\documentclass[../DoAn.tex]{subfiles}
\begin{document}

\begin{center}
    \Large{\textbf{TÓM TẮT NỘI DUNG ĐỒ ÁN}}\\
\end{center}
\vspace{1cm}

Đề tài "Xây dựng hệ thống hỗ trợ sắp xếp thời khoá biểu cho trường tiểu học" được thực hiện nhằm hỗ trợ giải quyết bài toán lập thời khóa biểu tại các trường tiểu học. Công việc này đòi hỏi người phụ trách phải phối hợp nhiều dữ liệu như lớp học, môn học, số tiết, giáo viên, phân công giảng dạy và phòng chức năng. Khi thực hiện bằng bảng tính, người phụ trách có thể chủ động điều chỉnh nhưng phải đối chiếu thủ công nhiều thông tin, làm tăng khối lượng kiểm tra và nguy cơ bỏ sót xung đột, đặc biệt khi giáo viên bộ môn giảng dạy nhiều lớp hoặc dữ liệu phân công thay đổi.

Qua khảo sát một số giải pháp hiện có, các hệ thống quản lý giáo dục tích hợp thường bao gồm phạm vi chức năng rộng, trong khi các phần mềm xếp lịch chuyên dụng được xây dựng cho bài toán tổng quát và có thể yêu cầu người dùng thiết lập nhiều dữ liệu, quy tắc trước khi áp dụng cho bậc tiểu học. Trên cơ sở đó, đề tài hướng đến xây dựng một ứng dụng web tập trung vào quy trình lập thời khóa biểu tiểu học, hỗ trợ quản lý dữ liệu đầu vào, phân công giảng dạy, sắp xếp thủ công trên giao diện lưới và tự động phân bổ các tiết học còn thiếu. Các vi phạm ràng buộc được hệ thống phát hiện và thông báo để người phụ trách kiểm tra, điều chỉnh và xác nhận trước khi công bố thời khóa biểu.

Về mặt công nghệ, hệ thống được xây dựng theo kiến trúc client–server tách biệt, sử dụng Next.js cho phía frontend và Spring Boot (Java) cho phía backend, xác thực qua JWT, cùng cơ sở dữ liệu quan hệ PostgreSQL. Việc lựa chọn kiến trúc này đảm bảo hệ thống dễ triển khai, dễ bảo trì và có thể vận hành hoàn toàn trên các nền tảng miễn phí, không phát sinh chi phí cho nhà trường.

Kết quả của đồ án là một hệ thống web tích hợp các chức năng quản lý dữ liệu, phân công giảng dạy, xây dựng, kiểm tra và công bố thời khóa biểu, phù hợp với các nghiệp vụ cốt lõi của trường tiểu học. Đồ án góp phần minh họa cách ứng dụng kiến trúc client--server và công cụ tối ưu hóa vào bài toán xếp thời khóa biểu, đồng thời tạo tiền đề cho việc tiếp tục hoàn thiện mô hình ràng buộc, mở rộng phạm vi thử nghiệm và phát triển các chức năng nâng cao trong tương lai.
\begin{flushright}
Sinh viên thực hiện\\
\begin{tabular}{@{}c@{}}
\textit{(Ký và ghi rõ họ tên)}
\end{tabular}
\end{flushright}

\end{document}
